\documentclass[11pt,a4paper]{article}
\usepackage[T1]{fontenc}
\usepackage{lmodern}
\usepackage[margin=2.6cm]{geometry}
\usepackage{amsmath,amssymb,amsthm,mathtools}
\usepackage{booktabs}
\usepackage{graphicx}
\usepackage{xcolor}
\usepackage{enumitem}
\usepackage{pgfplots}
\pgfplotsset{compat=1.17}
\usepackage[numbers,sort&compress]{natbib}
\usepackage{microtype}
\usepackage[hidelinks]{hyperref}

% numbers generated by theory/verify/*.py (see Appendix B); \input fails loudly if missing
\newcommand{\ExactAtFiveTwelve}{$2.44\times 10^{-11}$}
\newcommand{\ChernoffAtFiveTwelve}{$1.13\times 10^{-7}$}
\newcommand{\ExactAtTwoFiftySix}{$1.08\times 10^{-6}$}
\newcommand{\ExactAtSixtyFour}{$7.70\times 10^{-3}$}
\newcommand{\SimAtSixtyFour}{$7.71\times 10^{-3}$}
\newcommand{\GapAtFiveTwelve}{4\,637}
\newcommand{\GapAtSixtyFour}{315}
\newcommand{\AsymAtFiveTwelve}{$3.38\times 10^{-11}$}
\newcommand{\PoissonAtFiveTwelve}{$2.93\times 10^{-11}$}

\newcommand{\ExactTriples}{258}
\newcommand{\ExactMismatch}{0}

\newcommand{\UndoComparisons}{488\,919}
\newcommand{\UndoFailures}{0}
\newcommand{\ExactnessFailures}{0}
\newcommand{\ExactnessCases}{18}
\newcommand{\OrphanFailures}{0}

\newcommand{\LawCheckTrials}{400\,000}
\newcommand{\LawCheckMaxZ}{1.96}
\newcommand{\LawCheckMismatch}{0}

\newcommand{\HashSaveLo}{9.2}
\newcommand{\HashSaveHi}{12.2}
\newcommand{\PackSaveLo}{5.3}
\newcommand{\PackSaveHi}{9.0}
\newcommand{\TotSaveLo}{16.4}
\newcommand{\TotSaveHi}{17.9}
\newcommand{\InternLo}{-0.38}
\newcommand{\InternHi}{0.00}

\newcommand{\AdM}{2048}
\newcommand{\AdL}{512}
\newcommand{\AdHonest}{700}
\newcommand{\AdSeeds}{5}
\newcommand{\AdQueries}{2071}

\newcommand{\LedgerInstants}{848\,449}
\newcommand{\AsymRatioFour}{0.981}
\newcommand{\KnuthCases}{481}
\newcommand{\KnuthWorst}{2e-10}
\newcommand{\UniLHalf}{20}
\newcommand{\UniLThreeQ}{49}
\newcommand{\OrphanMeasured}{1.50}
\newcommand{\OrphanTheory}{1.50}


\newtheorem{theorem}{Theorem}
\newtheorem{proposition}[theorem]{Proposition}
\newtheorem{corollary}[theorem]{Corollary}
\theoremstyle{definition}
\newtheorem{definition}[theorem]{Definition}
\theoremstyle{remark}
\newtheorem{remark}{Remark}

\newcommand{\st}{STRATA}
\newcommand{\Z}{\mathbb{Z}}
\newcommand{\N}{\mathbb{N}}
\newcommand{\Prob}{\mathbb{P}}
\newcommand{\E}{\mathbb{E}}
\newcommand{\ceil}[1]{\lceil #1\rceil}
\newcommand{\code}[1]{\texttt{#1}}
\DeclareMathOperator{\Ins}{Ins}
\DeclareMathOperator{\Occ}{Occ}

\title{Stack-Disciplined Open Addressing:\\Exact Undo, Exact Refusal Probabilities, and Key-Lifetime Rules\\for Bounded-Memory Scoped Tables}
\author{Shubhi Singh \and Adhyan Jain\\[2pt]\small Vellore Institute of Technology, Vellore, India}
\date{}

\begin{document}
\maketitle

\begin{abstract}
\noindent Scoped symbol tables release declarations in strict last-in-first-out order. We ask what this discipline gives an open-addressing hash table that must live in one fixed buffer, must refuse a request cleanly when it is full, and must bound the work of every call. \emph{(1)} For every non-reordering open-addressing scheme, deletion in stack order by clearing one cell is an exact undo: the table is at all times identical to the insertion-only table of the live sequence, so insertion-only analyses of the table's state, and of the cost of the next operation, apply unchanged. Stack order is moreover \emph{exactly} the class of deletion orders for which clearing is safe for every hash assignment; leaving it orphans, in linear probing, on average $\alpha/(2(1-\alpha))$ live keys per deletion. \emph{(2)} For a table that refuses an insertion needing more than $L$ probes, the refusal probability is the tail of the classical unsuccessful-search cost of linear probing; we give an elementary derivation of its explicit finite-size form and use it as a design tool. The Chernoff-type bound used previously overstates it by a factor that grows like $L^{3/2}$; at load $3/4$ the limit needed for a per-insertion refusal probability of $10^{-6}$ is $258$ rather than $455$ for large tables, and far smaller for small ones. \emph{(3)} The guarantee of a keyed hash holds against adversarial names only if the key's lifetime does not exceed the table's: with a persistent key and an oracle that reveals home slots an adversary builds arbitrarily long runs with $O(m\log m)$ queries, and with only a probe-length oracle it learns one colliding name per $m$ queries. \emph{(4)} Exact space-ledger identities show that on eight compiler traces the $16$--$18\%$ saving of the compact layout over a conventional one is a stored hash word ($\HashSaveLo$--$\HashSaveHi\%$) plus name packing ($\PackSaveLo$--$\PackSaveHi\%$); interning itself is space-neutral. Results are checked by exhaustive enumeration, exact rational arithmetic, and a statistical test against a real implementation.
\end{abstract}

\noindent\textbf{Keywords:} open addressing; linear probing; stack discipline; symbol tables; bounded memory; hash flooding; exact analysis.

%%%%%%%%%%%%%%%%%%%%%%%%%%%%%%%%%%%%%%%%%%%%%%%%%%%%%%%%%%%%%%%%%%%%%%%%
\section{Introduction}\label{sec:intro}

A scoped symbol table binds identifiers to declarations and forgets the declarations of a block when the block ends. The textbook implementation is a hash table from names to their innermost declaration plus an undo stack that restores what a scope hid~\cite{appel1998,graham1979}. When the compiler or interpreter that owns the table runs on a microcontroller, in a boot loader, or under a fixed allocation, three further requirements appear: the table must live in \emph{one fixed buffer}, it must \emph{refuse} a request that does not fit and stay usable, and the \emph{work} of one call must be bounded even when the source text is supplied by someone else. \st~\cite{strata} is a design that meets them: one buffer holds an open-addressing index, one packed record per visible name, and the undo stack, and an operation that does not fit changes nothing.

An earlier, empirical study of \st\ reported what that design costs and saves on events logged by real compilers. This article asks a different question: \emph{what does the stack discipline itself imply?} Scope exit removes declarations in exactly the reverse order of their creation. That is a very special deletion order for an open-addressing table, and, as we show, it makes four things exact that are otherwise hard or open: the state of the table after deletions (Section~\ref{sec:undo}), the probability that a probe-limited table refuses an honest name (Section~\ref{sec:law}), what an adaptive adversary can do (Section~\ref{sec:keys}), and where the memory goes (Section~\ref{sec:ledger}).

\paragraph{Contributions.}
\begin{enumerate}[leftmargin=*,label=\textbf{C\arabic*.}]
\item \textbf{Exact undo and its converse} (Section~\ref{sec:undo}). For every non-reordering open-addressing scheme, popping the newest key by clearing its cell restores the previous table exactly (Theorem~\ref{thm:undo}); hence, in the random-hash model, the table at every moment has the law of an insertion-only table, and every insertion-only result about the table's state or the cost of the next operation transfers (Corollary~\ref{cor:marginal}). Stack order is the \emph{only} deletion order for which clearing is safe for all hash assignments (Theorem~\ref{thm:exact}), and clearing a random live key instead orphans on average the mean displacement of a key (Proposition~\ref{prop:orphans}). Combined with published results on open addressing without reordering we obtain a separation by deletion order (Corollary~\ref{cor:sep}).
\item \textbf{Exact refusal probabilities for probe-limited tables} (Section~\ref{sec:law}). For linear probing on a ring with $n$ keys among $m$ cells, the probability that an independent new key needs more than $L$ probes is the tail of the unsuccessful-search cost, whose exact generating function and Poisson limit are known~\cite{janson2016,knuth3}. We give an elementary, self-contained derivation of an explicit finite-size sum (Theorem~\ref{thm:law}) and its tail $\sim C\rho^{L}L^{-3/2}$ (Corollary~\ref{cor:limit}); what is new is the use: the Chernoff-type bound used for probe limits loses a factor $\Theta(L^{3/2})$, the exact values give smaller, tabulated limits (Table~\ref{tab:design}), and the formula is validated against a real probe-limited implementation.
\item \textbf{Key-lifetime rules} (Section~\ref{sec:keys}). With an independent key per table, the exact law holds against arbitrary adversarial names (Proposition~\ref{prop:fresh}); with a persistent key and a probe oracle it does not (Proposition~\ref{prop:persist}), and we quantify the cost of the attack.
\item \textbf{Exact space ledger} (Section~\ref{sec:ledger}). Identities for the words used by five single-buffer layouts at every instant (Proposition~\ref{prop:ledger}), evaluated on eight compiler traces; they decompose the saving of the compact layout and show that interning contributes none of it on these programs.
\end{enumerate}

\paragraph{What is not claimed.} The undo observation is elementary; its value lies in what follows from it and in the converse. The refusal probability itself is classical: it is the tail of the unsuccessful-search cost, whose exact finite-size generating function and Poisson limit are due to Knuth's run-length analysis~\cite{knuth3} and to Janson and Viola~\cite{janson2016}; we claim only an elementary derivation of an explicit finite-size sum, the quantitative gap to the Chernoff-type bound, the design rule, and the validation against an implementation. The adversary results are model-based (ideal hash, idealised oracle), not timing measurements. The proofs are pen-and-paper and are backed by exhaustive and statistical checks, not by a proof assistant. Section~\ref{sec:limits} lists what the evidence does not cover.

\paragraph{Outline.} Section~\ref{sec:model} fixes the model. Sections~\ref{sec:undo}--\ref{sec:ledger} give the four groups of results. Section~\ref{sec:related} places them in the literature, Section~\ref{sec:limits} states limitations, and Section~\ref{sec:conclusion} concludes. Appendix~\ref{app:proof} proves Theorem~\ref{thm:law}; Appendix~\ref{app:repro} explains how to regenerate every number.

%%%%%%%%%%%%%%%%%%%%%%%%%%%%%%%%%%%%%%%%%%%%%%%%%%%%%%%%%%%%%%%%%%%%%%%%
\section{Model}\label{sec:model}

\paragraph{Tables and schemes.} A table has $m\ge 2$ cells, indexed by $\Z_m$. Its state $T$ maps each cell to a name or to $\bot$ (empty); $\Occ(T)\subseteq\Z_m$ is the set of non-empty cells. Every name $x$ has a \emph{probe sequence} $\pi_x=(\pi_x(1),\pi_x(2),\dots)$ of cells, determined by the hash function (and its key). For linear probing, $\pi_x(i)=h(x)+i-1 \bmod m$ for a \emph{home} $h(x)\in\Z_m$.

\begin{definition}[Non-reordering scheme]\label{def:scheme}
An open-addressing scheme is \emph{non-reordering} if an insertion of $x$ into $T$ writes $x$ into one cell $\rho(\Occ(T),\pi_x)\in\{\pi_x(i)\}\setminus\Occ(T)$ --- or refuses, if the rule returns $\bot$ --- and changes no other cell. The rule $\rho$ may depend on the table only through the occupancy set (so control state that is a function of the occupancy, such as the number of occupied cells or per-region counters, is allowed).
\end{definition}

This covers linear, quadratic and double hashing, uniform probing, and non-greedy rules that decide from occupancy counts alone (which is how we read the funnel and elastic schemes of~\cite{fkk2024}; we did not re-derive this from their pseudocode). It excludes cuckoo and hopscotch hashing~\cite{pagh2004,herlihy2008} and Robin Hood hashing~\cite{celis1985}, which move keys.

\paragraph{Stack discipline.} The operations are $\textsc{push}(x)$, which inserts a name that is not live, and $\textsc{pop}$, which removes the most recently pushed live name by setting its cell to $\bot$. A scope mark is a position in the live sequence, and leaving the scope is the sequence of pops back to that position. A push that the rule refuses (buffer exhausted, probe limit exceeded) changes nothing and the name does not become live. The \emph{live sequence} is the list of live names in the order pushed. For a sequence $a=(a_1,\dots,a_t)$ of distinct names, $\Ins(a)$ denotes the table obtained by inserting $a_1,\dots,a_t$ in this order into the empty table.

\paragraph{Hash model.} In the \emph{random-function model}, the probe sequences of distinct names are independent and identically distributed; for linear probing the homes are independent and uniform on $\Z_m$. An operation sequence is \emph{oblivious} if it does not depend on the hash function. A keyed hash is modelled as a random function drawn from the key; Section~\ref{sec:keys} says what this assumption does and does not give against an adversary.

\paragraph{Mapping to \st.} In \st\ the index is a ring of $m$ words ($m$ a power of two) probed linearly; a name that hides a visible one needs no index slot, so the pushes of the model are the \emph{first} declarations of a name; shadowing declarations are recorded on the undo stack only. Exiting a scope pops the undo stack and, for each first declaration, sets one index slot to zero. In the growable mode a rebuild re-inserts the names oldest first, which re-establishes $\Ins$ of the live sequence for the larger index.

%%%%%%%%%%%%%%%%%%%%%%%%%%%%%%%%%%%%%%%%%%%%%%%%%%%%%%%%%%%%%%%%%%%%%%%%
\section{Exact Undo and Its Converse}\label{sec:undo}

\begin{theorem}[Undo]\label{thm:undo}
Let $\rho$ be a non-reordering rule. After any sequence of pushes and pops, with refused pushes leaving the state unchanged, the table equals $\Ins(a)$, where $a$ is the live sequence.
\end{theorem}
\begin{proof}
Induction on the number of operations. The empty table is $\Ins(())$. A refused push changes neither the table nor $a$. If a push of $x$ succeeds, the table before it is $\Ins(a)$ by hypothesis; the push writes $x$ into $\rho(\Occ(\Ins(a)),\pi_x)$, which is $\Ins(a\cdot x)$ by definition. For a pop, $a=u\cdot x$ and the table is $\Ins(u\cdot x)$, which is $\Ins(u)$ with $x$ additionally in the cell that the insertion of $x$ chose, because insertion changes exactly one cell and that cell was empty in $\Ins(u)$. Clearing it gives $\Ins(u)$.
\end{proof}

Two features of the proof matter. It does not use the lookup rule, so \emph{any} lookup procedure that is correct for insertion-only tables --- stopping at the first empty cell for greedy schemes, level-by-level searches for non-greedy ones --- stays correct under stack discipline. And it does not use randomness, so it also holds for adversarial homes.

\begin{corollary}[Marginals]\label{cor:marginal}
In the random-function model, for an oblivious operation sequence, at every moment up to the first refusal the law of the table is the law of the insertion-only table of the live sequence. In particular, for linear probing, Knuth's expected probe counts~\cite{knuth3}, the distribution of displacements~\cite{janson2005,viola2005}, and the refusal law of Section~\ref{sec:law} hold at every moment, for the number $n$ of live keys at that moment.
\end{corollary}

This concerns the \emph{state} of the table and the cost of the \emph{next} operation; cumulative or amortised costs over a whole schedule do not transfer. Successive moments are correlated (a name that is popped and pushed again returns to the same cell if the live sequence below it is unchanged), so Corollary~\ref{cor:marginal} is a statement about marginals; Section~\ref{sec:law} uses it with a union bound.

\subsection{Why stack order is exactly right}

For linear probing, a \emph{lookup} of $x$ scans from $h(x)$ and stops at $x$ or at an empty cell. Call a deletion order $\sigma$ of $n$ keys $x_1,\dots,x_n$ (inserted in this order into $\Z_m$, $2\le n\le m$) \emph{clear-safe for the homes $h$} if after every prefix of $\sigma$ has been cleared, the lookup of every remaining key still finds it.

\begin{theorem}[Exactness]\label{thm:exact}
$\sigma$ is clear-safe for every assignment of homes if and only if $\sigma=(x_n,x_{n-1},\dots,x_1)$, the reverse insertion order.
\end{theorem}
\begin{proof}
($\Leftarrow$) Let $y$ be the newest live key, in cell $c$. Every other live key $x$ is older; its lookup path, the cells from $h(x)$ up to its own cell, consisted of cells that were occupied when $x$ was inserted. Cell $c$ was empty then (the later insertion of $y$ put it there), so $c$ lies on no older key's path, and clearing it cannot stop any lookup early.

($\Rightarrow$) Suppose $\sigma$ is not the reverse order and let $t$ be the first step at which it deviates: the first $t-1$ deletions removed $x_n,\dots,x_{n-t+2}$, and step $t$ clears some $x_i$ with $i<z:=n-t+1$. Choose homes so that $x_1,\dots,x_{z-1}$ have distinct homes (possible since $z-1<m$), hence each sits in its own home cell, and give $x_z$ the home of $x_i$. Keys inserted after $x_z$ do not affect the placement of $x_z$. At step $t$ the table is $\Ins(x_1,\dots,x_z)$ and $x_z$ lies in a cell different from $x_i$'s cell, after at least one probe. Clearing $x_i$'s cell empties the first cell of $x_z$'s path, so the lookup of $x_z$ stops there and fails. Hence $\sigma$ is not clear-safe for this assignment.
\end{proof}

\begin{remark}\label{rem:general}
The proof uses only that a lookup follows the probe sequence of the key and stops at the first empty cell, and that the first probe of a key may be any cell. It therefore holds verbatim for every greedy probing scheme (linear, quadratic, double hashing, uniform probing). Besides linear probing, we checked it exhaustively for quadratic probing (all assignments of homes and all deletion orders, $m\in\{4,8\}$, $n\le4$).
\end{remark}

Theorem~\ref{thm:exact} explains why every deletion method for open addressing that supports arbitrary order needs more than clearing: tombstones, or backward shifting (Knuth's Algorithm~R)~\cite{knuth3,jimenez2018}, whose long-run behaviour is analysed in~\cite{bender2021}. It also explains why \st\ cannot simply swap in a relocating deletion: its undo entries store index slots, and relocation would invalidate them. The next proposition quantifies the damage of leaving the discipline.

\begin{proposition}[Orphans]\label{prop:orphans}
In a linear-probing table, let $\delta_y$ be the displacement of live key $y$ (the number of cells strictly between its home and its cell, cyclically) and, for a live key $x$, let $\mathrm{orph}(x)$ be the number of other live keys whose lookup fails after the cell of $x$ is cleared. Then $\sum_x \mathrm{orph}(x)=\sum_y \delta_y$. Consequently clearing a uniformly random live key orphans on average the mean displacement of a key; for $n=\alpha m$ keys with independent uniform homes this tends to $\alpha/(2(1-\alpha))$ (for example $1.5$ at $\alpha=3/4$). Clearing the newest key orphans none.
\end{proposition}
\begin{proof}
A key $y$ is orphaned by clearing $x$ iff the cell of $x$ lies on the path of $y$ strictly before $y$'s own cell. There are $\delta_y$ such cells and each is occupied by exactly one live key, so the pairs $(x,y)$ number $\sum_y\delta_y$. The total displacement does not depend on the insertion order (it equals the number of key--cell boundary crossings, a function of the home counts alone~\cite{knuth3}); its mean over independent uniform homes is $n(C_n-1)$, where $C_n$ is Knuth's expected cost of a successful search, and $C_n-1\to\tfrac12\bigl(\tfrac1{1-\alpha}-1\bigr)=\tfrac{\alpha}{2(1-\alpha)}$. The last claim is Theorem~\ref{thm:undo} (or the first half of Theorem~\ref{thm:exact}).
\end{proof}

\subsection{Consequence: a separation by deletion order}

The best insertion-only bounds for open addressing without reordering are due to Farach-Colton, Krapivin and Kuszmaul~\cite{fkk2024}: a greedy scheme (funnel hashing) with worst-case expected probe complexity $O(\log^2\delta^{-1})$ and amortised $O(\log\delta^{-1})$ for the insertions up to load $1-\delta$, and a non-greedy scheme (elastic hashing) with worst-case expected $O(\log\delta^{-1})$ and amortised $O(1)$. For insertions \emph{and} deletions over an unbounded horizon the same paper states that even linear probing and uniform probing resist analysis, and that the optimal amortised expected probe complexity is $\delta^{-\Omega(1)}$, citing Theorem~3 of~\cite{tinypointers2025}. We checked that statement in the primary source: Theorem~3 there lower-bounds the size of variable-size tiny pointers by $\Omega(\log\delta^{-1})$ bits at load $1-\delta$ with $\delta<1/4$, for a workload that alternates insertions of random elements with deletions of random present elements. An open-addressing scheme induces such pointers by encoding the probe index $i$ in $O(\log i)$ bits, so $\E[\log i]=\Omega(\log\delta^{-1})$ and, by Jensen, $\E[i]\ge\delta^{-\Omega(1)}$; this is the reduction of~\cite{fkk2024}, which we have not re-proved.

\begin{corollary}[Separation by deletion order]\label{cor:sep}
Suppose the hash functions are independent random functions and the operation sequence is oblivious. (i) For a scheme that satisfies Definition~\ref{def:scheme} (we read funnel and elastic hashing as doing so), every push under stack discipline at live count at most $(1-\delta)m$ has, by Theorem~\ref{thm:undo}, exactly the expected probe count it would have in the insertion-only table of the same live sequence; so, under the conditions on $\delta$ stated in~\cite{fkk2024}, for funnel hashing it is $O(\log^2\delta^{-1})$ and for elastic hashing the worst case is $O(\log\delta^{-1})$, \emph{whatever the push/pop schedule}. (ii) On a long workload that alternates insertions of random names with deletions of random live names at load $1-\delta$, $\delta<1/4$, no open-addressing scheme without reordering has expected probe index below $\delta^{-\Omega(1)}$~\cite{fkk2024,tinypointers2025}.
\end{corollary}

Part (ii) is quoted, part (i) is new only as a consequence of Theorem~\ref{thm:undo}; the corollary says that, at loads near one, the deletion \emph{order} moves the achievable probe complexity from polynomial to polylogarithmic in $\delta^{-1}$. It is a statement about expected probe counts; it does not make a claim about worst-case schedules at maximum load.

%%%%%%%%%%%%%%%%%%%%%%%%%%%%%%%%%%%%%%%%%%%%%%%%%%%%%%%%%%%%%%%%%%%%%%%%
\section{An Exact Refusal Law}\label{sec:law}

A bounded-work table fixes a \emph{probe limit} $L$: an insertion examines at most $L$ cells from its home and is refused (nothing changed) if all are occupied. Because a stored key never moves and clearing only empties cells, every stored key lies within $L$ cells of its home, so lookups also stop after $L$ cells; the work of a call is bounded by a constant independent of the contents. The question is how often an \emph{honest} key is refused.

By Corollary~\ref{cor:marginal}, the refusal probability at a moment with $n$ live keys is that of an insertion-only linear-probing table with $n$ keys: the probability that an independent new key with uniform home finds its first $L$ cells occupied. Equivalently, it is the probability that the cost of an unsuccessful search exceeds $L$, a quantity whose distribution has been studied extensively~\cite{knuth3,fpv1998,janson2001}. Janson and Viola give the exact generating function for finite $m,n$ and, in the infinite Poisson model, $\Prob(U=k)=(1-\alpha)\Prob(B>k)$ for the block length $B$~\cite[Thm.~11.1, Eq.~(11.4)]{janson2016}; Knuth's run-length probabilities~\cite{knuth3} give the same finite-size values by Abel-binomial counting.

\paragraph{Notation.} For $g\ge1$ let $B(g)=g^{g-1}/g!$, and for $j\ge1$, $s\ge j$ let $Q_j(s)=\tfrac{j}{s}\,s^{s-j}/(s-j)!$ and $Q_0(s)=\mathbf 1[s=0]$. By the Borel--Tanner identity~\cite{tanner1961}, $Q_j(s)=\sum_{g_1+\dots+g_j=s}\prod_i B(g_i)$.

\begin{theorem}[Exact refusal law]\label{thm:law}
Let $m\ge2$ cells in a ring hold $n=m-k$ keys, $k\ge1$, with independent uniform homes under linear probing, and let a new key have an independent uniform home. The probability $R_m(n,L)$ that its first $L$ cells are all occupied is
\begin{equation}\label{eq:law}
R_m(n,L)\;=\;\frac{k}{m}\sum_{g=L+1}^{n+1}(g-L)\,\pi_g,\qquad
\pi_g=\frac{B(g)\,Q_{k-1}(m-g)}{Q_k(m)} .
\end{equation}
Here $\pi_g=\Prob(\gamma_1=g\mid\gamma_1+\dots+\gamma_k=m)$ for i.i.d.\ Borel variables $\gamma_i$, and does not depend on the Borel parameter. For $k=1$, $R_m(m-1,L)=(m-L)^+/m$.
\end{theorem}

\begin{remark}\label{rem:classical}
Theorem~\ref{thm:law} computes a known quantity: $R_m(n,L)$ is the tail of the unsuccessful-search cost, and Corollary~\ref{cor:limit}(b) is the form $(1-\alpha)\E[(B-L)^+]$ of~\cite[Eq.~(11.4)]{janson2016} with $B$ Borel distributed. What the proof below adds is a short derivation that uses only the cycle lemma and the Borel--Tanner identity and yields the explicit sum~\eqref{eq:law}; it was cross-checked against enumeration and against the implementation (Section~\ref{sec:law}).
\end{remark}

The proof (Appendix~\ref{app:proof}) cuts the ring at its $k$ empty cells; each empty cell and the run of occupied cells after it form a block whose size has the Borel distribution (the total progeny of a Poisson branching process, by the cycle lemma~\cite{dwass1969}), and the blocks, conditioned on filling the ring, are exchangeable.

\begin{corollary}[Limit, monotonicity, asymptotics]\label{cor:limit}
(a) $R_m(n,L)$ is nondecreasing in $n$, so the value at the maximum load bounds every lower load. (b) If $n/m\to\alpha\in(0,1)$ then $R_m(n,L)\to R_\infty(\alpha,L)=(1-\alpha)\sum_{g>L}(g-L)\,\mathrm{Borel}(g;\alpha)$, with $\mathrm{Borel}(g;\alpha)=e^{-\alpha g}(\alpha g)^{g-1}/g!$. (c) With $\rho=\alpha e^{1-\alpha}<1$,
\[
R_\infty(\alpha,L)\;\sim\;\frac{(1-\alpha)\,\rho^{L+1}}{\alpha\sqrt{2\pi}\,(1-\rho)^{2}\,L^{3/2}}\qquad(L\to\infty).
\]
\end{corollary}
\begin{proof}
(a) The occupied set of linear probing does not depend on the insertion order~\cite{konheim1966}. Insert the $(n+1)$-st key last: the occupied set with $n+1$ keys is the occupied set with $n$ keys plus one cell, and the event ``the first $L$ cells are occupied'' is increasing in the occupied set. (b) Write $\pi_g$ via the explicit $Q_j$; for fixed $g$, with $n=\alpha m$, $\pi_g\to B(g)\alpha^{g-1}e^{-\alpha g}=\mathrm{Borel}(g;\alpha)$. The $\pi_g$ are probability masses and $\E[\gamma_1\mid\sum\gamma_i=m]=m/k\to1/(1-\alpha)=\E[\mathrm{Borel}(\alpha)]$; pointwise convergence of masses together with convergence of the means gives uniform integrability (a standard consequence of Scheffé's lemma), hence convergence of $\E[(\gamma-L)^+]$, and $k/m\to1-\alpha$. (c) Stirling's formula gives $\mathrm{Borel}(g;\alpha)\sim\rho^g/(\alpha\sqrt{2\pi}\,g^{3/2})$; then $\sum_{j\ge1}j\,\mathrm{Borel}(L+j;\alpha)\sim \rho^L(\alpha\sqrt{2\pi}L^{3/2})^{-1}\sum_j j\rho^j$ by dominated convergence.
\end{proof}

\paragraph{Comparison with the Chernoff-type bound.} The repository's earlier analysis of the probe limit~\cite{strata} used the interval argument of linear probing with a Chernoff bound and obtained $\Prob(\text{refusal})\le\rho^L/(1-\rho)$. The bound is valid. Corollary~\ref{cor:limit}(c) shows that it overstates the truth by a factor $\sim\alpha\sqrt{2\pi}\,(1-\rho)L^{3/2}/((1-\alpha)\rho)\to\infty$: a union over all run lengths $\ge L$ pays a factor $1/(1-\rho)$ for events that overlap heavily, and the Chernoff estimate omits the polynomial factor $L^{-3/2}$ of the Borel tail. At load $3/4$ and $m=65\,536$ the exact value at $L=512$ is \ExactAtFiveTwelve{} against the bound's \ChernoffAtFiveTwelve{}, a factor of \GapAtFiveTwelve{}; at $L=64$ the factor is \GapAtSixtyFour{} (Figure~\ref{fig:law}).

\begin{figure}[t]
\centering
\begin{tikzpicture}
\begin{axis}[width=0.86\linewidth,height=6.4cm,ymode=log,xmin=8,xmax=600,ymin=1e-16,ymax=20,
  xlabel={probe limit $L$},ylabel={refusal probability},legend pos=south west,legend cell align=left,
  legend style={font=\small},grid=both,grid style={gray!20}]
\addplot[thick,blue] table[col sep=comma,x=L,y=exact]{data/fig_law.csv};
\addlegendentry{exact law, eq.~\eqref{eq:law}}
\addplot[thick,red,dashed] table[col sep=comma,x=L,y=chernoff]{data/fig_law.csv};
\addlegendentry{Chernoff-type bound $\rho^L/(1-\rho)$}
\addplot[only marks,mark=*,black,mark size=1.8pt] table[col sep=comma,x=L,y=simulated]{data/fig_sim.csv};
\addlegendentry{simulation, $10^7$ samples}
\end{axis}
\end{tikzpicture}
\caption{Probability that an insertion needs more than $L$ probes, $m=65\,536$ cells, load $3/4$. The exact law agrees with the stored simulation of the repository where the latter observes events; the Chernoff-type bound is looser by a factor that grows with $L$.}\label{fig:law}
\end{figure}

\paragraph{Design rule.} The smallest limit that keeps the per-insertion refusal probability below $\varepsilon$ is read off the exact law (Table~\ref{tab:design}, Figure~\ref{fig:design}). Two points are worth stating. For large tables the exact limit is about $57\%$ of the Chernoff-derived one at load $3/4$ ($258$ against $455$ for $\varepsilon=10^{-6}$). And small tables need much less: the finite-size law is smaller than its Poisson limit, so a $256$-cell index needs only $L=121$. \paragraph{Other probing rules.} For uniform probing, where the probe sequence of a name is a uniformly random sequence of distinct cells independent of the table, the probability that the first $L$ probes all hit occupied cells is exactly $\binom nL/\binom mL=\prod_{i<L}\frac{n-i}{m-i}\le\alpha^L$ \emph{whatever the occupied set}; by Theorem~\ref{thm:undo} it therefore holds at every moment under stack discipline. For $\varepsilon=10^{-6}$ and $m=2^{16}$ it gives $L=\UniLHalf$ at load $1/2$ and $L=\UniLThreeQ$ at load $3/4$, against $52$ and $258$ for linear probing (about one fifth of the work per call), at the price of the cache locality that makes linear probing attractive; double hashing and quadratic probing approximate uniform probing but are not covered by this exact value. A caller who wants a failure probability over a whole compilation of $N$ insertions uses $\varepsilon=\varepsilon'/N$, by the following union bound.

\begin{proposition}[Run of a compilation]\label{prop:run}
Under the random-function model with an oblivious sequence of $N$ pushes of names distinct from the live ones, and at most $n_{\max}$ live keys, the probability that some push is refused is at most $N\,R_m(n_{\max},L)$.
\end{proposition}
\begin{proof}
Until the first refusal the process coincides with the unlimited process. At each push the new name has an independent uniform home, independent of the live names' homes, and the table is the insertion-only table of the live sequence (Corollary~\ref{cor:marginal}); the refusal probability is $R_m(n_t,L)\le R_m(n_{\max},L)$ by Corollary~\ref{cor:limit}(a). Union bound over the pushes.
\end{proof}

\begin{table}[t]
\centering\small
\caption{Smallest probe limit $L$ with per-insertion refusal probability at most $\varepsilon$: exact finite-size law for index sizes $m$, its Poisson limit ($m\to\infty$), and the limit that the Chernoff-type bound would require.}\label{tab:design}
\begin{tabular}{@{}llrrrrrrr@{}}
\toprule
load & $\varepsilon$ & $m{=}2^{8}$ & $2^{10}$ & $2^{12}$ & $2^{14}$ & $2^{16}$ & $m\to\infty$ & Chernoff \\
\midrule
$3/8$ & $10^{-6}$ & 25 & 28 & 29 & 29 & 29 & 29 & 43 \\
$3/8$ & $10^{-9}$ & 37 & 44 & 46 & 47 & 47 & 47 & 62 \\
$1/2$ & $10^{-6}$ & 42 & 49 & 51 & 52 & 52 & 52 & 81 \\
$1/2$ & $10^{-9}$ & 59 & 77 & 83 & 84 & 85 & 85 & 117 \\
$5/8$ & $10^{-6}$ & 70 & 93 & 101 & 104 & 104 & 104 & 171 \\
$5/8$ & $10^{-9}$ & 95 & 142 & 162 & 168 & 170 & 170 & 244 \\
$3/4$ & $10^{-6}$ & 121 & 201 & 242 & 255 & 258 & 260 & 455 \\
$3/4$ & $10^{-9}$ & 148 & 289 & 380 & 413 & 422 & 426 & 638 \\
\bottomrule
\end{tabular}

\end{table}

\begin{figure}[t]
\centering
\begin{tikzpicture}
\begin{axis}[width=0.86\linewidth,height=5.6cm,ymode=log,xmin=0.2,xmax=0.85,xlabel={load $\alpha$},ylabel={smallest $L$ for $\varepsilon=10^{-6}$},
  legend pos=north west,legend cell align=left,legend style={font=\small},grid=both,grid style={gray!20}]
\addplot[thick,blue,mark=*,mark size=1.4pt] table[col sep=comma,x=load,y=finite4096]{data/fig_design.csv};
\addlegendentry{exact, $m=4096$}
\addplot[thick,teal,mark=square*,mark size=1.4pt] table[col sep=comma,x=load,y=poisson]{data/fig_design.csv};
\addlegendentry{exact, $m{\to}\infty$}
\addplot[thick,red,dashed,mark=triangle*,mark size=1.6pt] table[col sep=comma,x=load,y=chernoff]{data/fig_design.csv};
\addlegendentry{Chernoff-type bound}
\end{axis}
\end{tikzpicture}
\caption{Probe limit required for a per-insertion refusal probability of $10^{-6}$.}\label{fig:design}
\end{figure}

\subsection{How the law was checked}

Four independent checks support Theorem~\ref{thm:law}. (i) \emph{Exhaustive enumeration:} for every ring size $m\le 9$ and every $n,L$ with $m^n\le2\cdot10^6$ homes sequences ($\ExactTriples$ triples $(m,n,L)$), formula~\eqref{eq:law} evaluated in exact rational arithmetic equals the probability obtained by simulating linear probing on \emph{every} sequence of homes: $\ExactMismatch$ mismatches. (ii) \emph{The stored simulation:} the repository's $240$-million-sample simulation of fresh insertions into tables of $65\,536$ slots agrees with the exact law wherever the simulation observed events (e.g.\ $7.71\times10^{-3}$ against \ExactAtSixtyFour{} at $L=64$, load $3/4$). (iii) \emph{The real implementation:} Table~\ref{tab:lawcheck} compares the law with the actual \st\ table (keyed hash with a random key per trial, static index, probe limit): the observed frequency of a fresh name needing more than $L$ probes in \LawCheckTrials{} independent trials per row. The largest $|z|$ is \LawCheckMaxZ. In the same trials a second table with the probe limit set and an identical key was kept in lockstep: while no earlier probe exceeded $L$ its state digest equalled that of the unlimited table, and it returned the refusal status exactly when the probe count exceeded $L$ (\LawCheckMismatch{} mismatches). (iv) \emph{An independent derivation:} counting by Abel's binomial theorem in the manner of Knuth~\cite{knuth3} gives the probability that a run of exactly $g$ occupied cells starts at a fixed cell as $\binom ng(g+1)^{g-1}(k-1)(m-g-1)^{n-g-1}/m^{n}$, hence $\sum_{g\ge L}(g-L+1)P_g$ for the refusal probability; this shares no step with the proof in Appendix~\ref{app:proof} and agrees with formula~\eqref{eq:law} in \KnuthCases{} cases (exactly for $m\le14$, to relative difference \KnuthWorst{} for $m$ up to $2^{20}$).

\begin{table}[t]
\centering\small
\caption{Exact law against the real table. $m$: index slots, $n$: keys already stored, $L$: probe limit. ``Clean'' trials are those in which no earlier insertion exceeded $L$ (used for the lockstep check).}\label{tab:lawcheck}
\begin{tabular}{@{}rrrrrrrrr@{}}
\toprule
$m$ & $n$ & $L$ & trials & observed & exact & $z$ & clean trials & mismatches \\
\midrule
256 & 128 & 8 & 400\,000 & 0.03311 & 0.03315 & -0.13 & 254\,871 & 0 \\
256 & 160 & 16 & 400\,000 & 0.02643 & 0.02667 & -0.94 & 307\,410 & 0 \\
256 & 190 & 24 & 400\,000 & 0.05740 & 0.05730 & +0.26 & 247\,303 & 0 \\
256 & 190 & 40 & 400\,000 & 0.01372 & 0.01375 & -0.19 & 369\,939 & 0 \\
1024 & 760 & 48 & 400\,000 & 0.01251 & 0.01286 & -1.96 & 307\,286 & 0 \\
1024 & 760 & 96 & 400\,000 & 0.00060 & 0.00058 & +0.52 & 397\,066 & 0 \\
1024 & 512 & 24 & 400\,000 & 0.00040 & 0.00046 & -1.63 & 395\,607 & 0 \\
4096 & 3071 & 80 & 400\,000 & 0.00317 & 0.00302 & +1.70 & 335\,646 & 0 \\
4096 & 3071 & 100 & 400\,000 & 0.00110 & 0.00106 & +0.60 & 378\,922 & 0 \\
16384 & 12287 & 80 & 100\,000 & 0.00296 & 0.00328 & -1.79 & 46\,819 & 0 \\
\bottomrule
\end{tabular}

\end{table}

%%%%%%%%%%%%%%%%%%%%%%%%%%%%%%%%%%%%%%%%%%%%%%%%%%%%%%%%%%%%%%%%%%%%%%%%
\section{Keys and Adversaries}\label{sec:keys}

Theorem~\ref{thm:law} assumes independent uniform homes. Against names chosen by an adversary this is a statement about the hash, and a statement about \emph{what the adversary may observe}. We make this precise.

\paragraph{Adversary model.} The table's hash is a random function $F_K$ drawn from a secret key $K$ (a pseudorandom function is modelled as ideal). An adversary chooses the names that are declared before the honest ones. A \emph{strong oracle} returns $F_K(\text{name})$ for a queried name; a \emph{weak oracle} returns one bit, whether the home of the queried name equals the home of a fixed anchor name (this is what a perfect timing channel on a probe-limited table reveals: the declaration is slower exactly when the probe sequence is longer).

\begin{proposition}[Fresh key]\label{prop:fresh}
Let every table be created with a key drawn independently of everything the adversary has seen, and let the adversary fix its names before the key is drawn. Then the homes of all distinct names, honest or not, are independent and uniform, and Theorem~\ref{thm:law} and Proposition~\ref{prop:run} hold verbatim: whatever names the adversary declares and however it schedules scopes, the refusal probability of an honest name is at most $R_m(n_{\max},L)$.
\end{proposition}
\begin{proof}
In the ideal model $F_K$ for an independent $K$ is a uniformly random function; an oblivious choice of names cannot correlate with it. The remaining statements are Theorem~\ref{thm:law} and Corollary~\ref{cor:limit}(a) applied to the live count, which the adversary can only increase up to $n_{\max}$.
\end{proof}

\begin{proposition}[Persistent key, strong oracle]\label{prop:persist}
Suppose the same key is used for many tables and the adversary has a strong oracle. After $q\ge m(\ln m+t)$ queries of random names, with probability at least $1-e^{-t}$ the adversary knows a name for every cell. Declaring names for $r\le n_{\max}$ consecutive cells first, in an otherwise empty table, produces a run of $r$ occupied cells without any refusal, and then the next honest name is refused with probability at least $(r-L+1)/m$ when $r\ge L$. For $r$ well above $L$ this is far above the exact law (Theorem~\ref{thm:law}).
\end{proposition}
\begin{proof}
A fixed cell is missed by $q$ independent uniform queries with probability $(1-1/m)^q\le e^{-q/m}$; a union over $m$ cells gives $me^{-q/m}\le e^{-t}$. Names for $r$ consecutive cells each go to their own home, so the run occupies exactly those cells. An honest name whose home is one of the first $r-L+1$ cells of the run finds $L$ occupied cells.
\end{proof}

\begin{proposition}[Weak oracle]\label{prop:weak}
With a weak oracle each query reveals a collider of the anchor with probability $1/m$ independently, so finding $c$ colliders takes $cm$ queries in expectation. A probe-limited table can store at most $L-1$ of them at the anchor's home, because the $j$-th needs $j+1$ probes.
\end{proposition}

\paragraph{Illustration.} Table~\ref{tab:adaptive} runs the weak-oracle attack against the real \st\ table with a secret persistent key ($m=\AdM$, $L=\AdL$, \AdHonest{} honest names after the adversary's names, \AdSeeds{} seeds). Learning one collider cost on average \AdQueries{} queries, as the geometric argument predicts ($\approx m$). With the same key, between $1\%$ and $5.5\%$ of the honest names are refused even though memory is available; names the attacker learned are useless against a table with a fresh key, and the same number of random names causes no refusal. The oracle here is the library's \code{longest\_probe} statistic, an idealised stand-in for a timing channel; this is an illustration of Proposition~\ref{prop:weak}, not a timing attack.

\begin{table}[t]
\centering\small
\caption{Honest names refused by the probe limit (percent of \AdHonest): the adversary's learned colliders declared first. Mean over \AdSeeds{} seeds; maximum over seeds for the persistent key.}\label{tab:adaptive}
\begin{tabular}{@{}rrrrr@{}}
\toprule
colliders & persistent key (mean) & (max) & fresh key & random names \\
\midrule
400 & 1.00\% & 1.29\% & 0.00\% & 0.00\% \\
512 & 5.51\% & 6.14\% & 0.00\% & 0.00\% \\
600 & 4.71\% & 6.00\% & 0.00\% & 0.00\% \\
700 & 5.14\% & 5.71\% & 0.00\% & 0.00\% \\
800 & 5.14\% & 6.00\% & 0.00\% & 0.00\% \\
\bottomrule
\end{tabular}

\end{table}

\paragraph{Key-lifetime rule.} Propositions~\ref{prop:fresh}--\ref{prop:persist} say the guarantee of the probe limit against adversarial \emph{names} requires that no information about a key's behaviour leaks into names that are declared in a table that uses the same key. In practice: derive the key of the $i$-th table from a device secret and a never-repeating counter, $K_i=\mathrm{PRF}(K_0,i)$, and never reuse a counter value (a device that resets its counter reuses keys). What does not depend on the key is the \emph{work bound}: a call examines at most $L$ cells whatever the names, so the worst case of hash flooding~\cite{crosby2003,aumasson2012} is bounded even when the key leaks; only the absence of refusals depends on it. Adaptivity against hash tables through timing is an old observation~\cite{lipton1993,naor2015}; the point here is the exact accounting of what it costs against a probe-limited table.

%%%%%%%%%%%%%%%%%%%%%%%%%%%%%%%%%%%%%%%%%%%%%%%%%%%%%%%%%%%%%%%%%%%%%%%%
\section{An Exact Space Ledger}\label{sec:ledger}

The layouts compared are single-buffer scoped tables whose size at an instant is a sum of word counts. For a name of $\ell$ bytes let $w_4(\ell)=\ceil{\ell/4}$ (raw words) and $w_5(\ell)=\ceil{\ell/5}$ (five six-bit characters per word, for names of letters, digits and underscore). At an instant let $d$ be the number of live declarations, $n$ the number of distinct visible names, $h=d-n$ the hidden declarations, $s$ the open scopes, $m$ the index slots (a power of two, at least $16$, doubled when $4(n+1)>3m$, never reduced within a translation unit) and $b$ the buckets of a chained table (doubled when $d+1>b$). The five layouts are:

\begin{center}\footnotesize
\setlength{\tabcolsep}{4pt}
\begin{tabular}{@{}lll@{}}
\toprule
layout & words in use & per-element costs\\
\midrule
$C_0$ chained, stored hash (baseline) & $b+1+\sum\nolimits_{\text{decl}}(4+w_4)+s$ & link, hash, value, meta, name\\
$C_1$ chained, no stored hash & $b+1+\sum\nolimits_{\text{decl}}(3+w_4)+s$ & link, value, meta, name\\
$C_2$ open-addressed, entry per declaration & $m+1+\sum\nolimits_{\text{decl}}(3+w_4)+s$ & prev, value, meta, name\\
$S_{\mathrm{raw}}$ interned, raw names & $m+1+\sum\nolimits_{x}(2+w_4)+n+3h+2s$ & record per name, undo\\
$S$ interned, packed names (\st) & $m+1+\sum\nolimits_{x}(2+w)+n+3h+2s$ & $w=w_5$ if packable else $w_4$\\
\bottomrule
\end{tabular}
\end{center}

Only $C_0$, $S_{\mathrm{raw}}$ and $S$ exist as code in the repository; $C_1$ and $C_2$ are layout rules, evaluated by accounting alone.

\begin{proposition}[Ledger]\label{prop:ledger}
At every instant, with the same $m$ and $b$ as the layouts would have at that instant,
\begin{align*}
U_{C_0}-U_{C_1}&=d, & U_{C_1}-U_{C_2}&=b-m,\\
U_{C_2}-U_{S_{\mathrm{raw}}}&=\sum_{\text{hidden decl.}}w_4(\ell)\;-\;s, &
U_{S_{\mathrm{raw}}}-U_{S}&=\sum_{x\text{ visible, packable}}\bigl(w_4(\ell_x)-w_5(\ell_x)\bigr).
\end{align*}
\end{proposition}
\begin{proof}
Direct subtraction of the formulas above: $C_0$ and $C_1$ differ by one word per declaration; $S_{\mathrm{raw}}$ charges $3+w_4$ per first declaration (record $2+w_4$ plus one undo word) and $3$ per hidden declaration, where $C_2$ charges $3+w_4$ for every declaration, and one scope mark costs $2$ words in $S$ against $1$ in $C_2$; packing changes $w_4$ to $w_5$ for packable names only.
\end{proof}

Hence the saving of the compact layout over a conventional one has an exact anatomy: a stored hash costs exactly one word per declaration; interning with an undo log pays only through the \emph{names of hidden declarations} and costs one extra word per open scope; and packing saves $\ceil{\ell/4}-\ceil{\ell/5}$ words per name, about $\ell/20$.

\begin{table}[t]
\centering\small
\caption{Smallest fixed buffer (KiB) for the five layouts on the eight compiler traces of the repository, and the savings (percent): of dropping the stored hash ($C_0\to C_1$), of packing ($C_1\to S$), and in total ($C_0\to S$). The accounting reproduces all $24$ minimum buffers measured for $C_0$, $S_{\mathrm{raw}}$ and $S$ exactly (to the $64$-byte resolution of the measurement).}\label{tab:ladder}
\begin{tabular}{@{}lrrrrrrrr@{}}
\toprule
trace & $C_0$ & $C_1$ & $C_2$ & $S_{\mathrm{raw}}$ & $S$ & hash & pack & total \\
\midrule
zlib & 18.9 & 16.6 & 16.6 & 16.7 & 15.8 & 11.9 & 5.3 & 16.6 \\
cJSON & 35.9 & 31.8 & 31.8 & 31.8 & 30.0 & 11.5 & 5.5 & 16.4 \\
littlefs & 24.2 & 21.4 & 21.4 & 21.5 & 20.0 & 11.6 & 6.7 & 17.5 \\
FreeRTOS & 53.6 & 47.8 & 47.8 & 47.8 & 44.0 & 11.0 & 7.9 & 17.9 \\
Lua & 43.4 & 38.1 & 38.1 & 38.1 & 35.9 & 12.2 & 5.7 & 17.3 \\
Mbed TLS & 96.9 & 87.9 & 87.9 & 87.9 & 80.0 & 9.2 & 9.0 & 17.4 \\
TinyCC & 189.8 & 168.2 & 168.2 & 168.3 & 156.5 & 11.3 & 7.0 & 17.5 \\
SQLite & 169.4 & 152.8 & 152.8 & 152.9 & 140.4 & 9.8 & 8.1 & 17.2 \\
\bottomrule
\end{tabular}

\end{table}

Table~\ref{tab:ladder} evaluates the ledger on the traces. The compact layout needs $\TotSaveLo$--$\TotSaveHi\%$ less than the conventional chained table, of which the stored hash word accounts for $\HashSaveLo$--$\HashSaveHi\%$ and packing for $\PackSaveLo$--$\PackSaveHi\%$ (measured against $C_1$). Interning with an undo log changes the footprint by between $\InternHi\%$ and $\InternLo\%$ relative to the open-addressed layout $C_2$ (negative: larger, because a scope mark costs two words): the programs shadow little, and Proposition~\ref{prop:ledger} says why --- the interned layout saves $\sum_{\text{hidden}}w_4-s$ words. This corrects an impression that the earlier empirical comparison could leave: against a conventional table that stores a hash, the compact layout looks one sixth smaller, but that margin is the stored hash plus packing, not the interning or the undo entries. The ledger does not say what happens on programs with heavy shadowing; the identity does: the saving then grows linearly in the number and length of hidden names.

%%%%%%%%%%%%%%%%%%%%%%%%%%%%%%%%%%%%%%%%%%%%%%%%%%%%%%%%%%%%%%%%%%%%%%%%
\section{Related Work}\label{sec:related}

\paragraph{Scoped symbol tables.} The undo-stack scheme is described by Appel~\cite{appel1998} and goes back to hashed tables with explicit scope control~\cite{graham1979}; LLVM's \code{ScopedHashTable}~\cite{llvmsht} implements it over an allocator. We read the symbol table of lcc~\cite{fraserhanson1995,lccsym}, the closest arena-based design: names are interned strings hashed by address, a table is created lazily per scope level and linked to the enclosing one by a \code{previous} pointer, lookup walks that chain outward, and leaving a scope pops the table pointer; storage comes from arenas~\cite{hanson1990} released at function end. It has neither a single index nor an undo log, a lookup may have to search every enclosing scope table, and it has no budget or refusal semantics.

\paragraph{Linear probing and deletions.} The analysis of linear probing~\cite{knuth3,pittel1987,fpv1998,janson2001,janson2005,viola2005} gives the cost distributions for insertion-only tables; our Theorem~\ref{thm:law} is the tail of the unsuccessful-search cost, whose exact generating function and Poisson limit are in~\cite{janson2016}; we contribute an elementary derivation of the explicit finite-size sum and its use as a refusal probability, not the distribution itself. For deletions, Knuth's backward shifting and tombstones~\cite{knuth3,jimenez2018}, the behaviour of tombstones~\cite{bender2021}, and tombstone minimisation with stable positions~\cite{sanders2018} address arbitrary orders; the state of the art for open addressing without reordering is~\cite{fkk2024} (insertions) and~\cite{tinypointers2025} (the lower bound for dynamic workloads), with a follow-up on greedy open addressing in 2026~\cite{fkk2026}. Dynamic linear probing has been analysed under random deletions in the uniform-hashing model~\cite{larson1983}. None of these treats the stack-ordered case; a few web searches (``LIFO deletion open addressing'', ``undo/rollback hash table'') found no statement of Theorems~\ref{thm:undo}--\ref{thm:exact}, which is weaker evidence than a systematic search and may reflect folklore in backtracking implementations. History-independent hashing~\cite{naor2001,blelloch2007} makes the layout a function of the stored \emph{set}; the canonical state of Theorem~\ref{thm:undo} is a function of the stored \emph{sequence}, which is weaker.

\paragraph{Bounded probe sequences.} Alharbi, Banderier and Bouillaguet~\cite{abb2026} study a bounded-probe linear-probing \emph{process} in which a key whose probe sequence exceeds $p$ is discarded and another is drawn, and derive the expected number of attempts needed to fill the table; they list the probability that an insertion succeeds at a given fill ratio as a question left open. That is a different quantity from ours: Theorem~\ref{thm:law} gives the refusal probability of an insertion into the table that the unlimited process would have built, which equals the refusal probability of the bounded process up to its first refusal (Proposition~\ref{prop:run}); after a refusal the two processes differ and we make no claim. Cuckoo~\cite{pagh2004} and hopscotch~\cite{herlihy2008} hashing bound probes by moving keys, which would break the undo property. A patent describes a cache with probe-limited open addressing that evicts or refuses an insertion at the limit~\cite{nemes2015}; to our knowledge no analysis of the refusal probability of such tables has been published.

\paragraph{Hash flooding and adaptivity.} Algorithmic complexity attacks~\cite{crosby2003} and seed-independent multicollisions of non-cryptographic hashes~\cite{aumasson2012} motivate keyed hashing; timing as an oracle for collisions~\cite{lipton1993} and adversarial Bloom filters~\cite{naor2015} are the adaptive-adversary precedents. Our contribution here is the accounting for a probe-limited table.

%%%%%%%%%%%%%%%%%%%%%%%%%%%%%%%%%%%%%%%%%%%%%%%%%%%%%%%%%%%%%%%%%%%%%%%%
\section{Limitations}\label{sec:limits}

\begin{itemize}[leftmargin=*]
\item \textbf{Ideal hashing.} Theorem~\ref{thm:law} and Corollary~\ref{cor:marginal} assume independent uniform homes. Real identifiers are not random strings; the repository's compiler traces are consistent with uniform spreading but do not prove it. A keyed hash is treated as an ideal function; HalfSipHash-1-3 with a 64-bit key is documented as adequate for hash-table flooding only.
\item \textbf{Marginals, not trajectories.} The law gives the refusal probability at one moment; the run bound (Proposition~\ref{prop:run}) is a union bound and ignores the correlation between moments. A sharper statement about the number of refusals in a run is not derived.
\item \textbf{The law is a known quantity.} Theorem~\ref{thm:law} and Corollary~\ref{cor:limit}(b) restate classical results (Remark~\ref{rem:classical}); the new parts are the Chernoff gap, the design rule, the validation against an implementation, and the link to stack discipline (Corollary~\ref{cor:marginal}).
\item \textbf{Two probing rules in the refusal law.} Theorems~\ref{thm:undo} and~\ref{thm:exact} cover the general class and greedy schemes respectively; the refusal probability is exact for linear probing (Theorem~\ref{thm:law}) and for uniform probing (a one-line product), not for quadratic probing or double hashing.
\item \textbf{The separation is conditional.} Corollary~\ref{cor:sep}(ii) rests on published results whose reduction we did not re-prove; the comparison is of expected probe counts, in the random-hash model.
\item \textbf{Idealised oracle.} The attacks of Section~\ref{sec:keys} use a model oracle. No timing measurement was made, and whether a given device leaks enough timing to realise it is not studied.
\item \textbf{Ledger.} The five layouts are accounting rules; $C_1$ and $C_2$ were not implemented, so their speed is unknown. The conclusions about interning hold for the eight C traces and the stated identity; heavily shadowing code is not covered by the data.
\item \textbf{Verification.} The proofs are on paper. Exhaustive enumeration ($m\le9$), randomized checks and the statistical test against the implementation support them; none replaces a machine-checked proof.
\item \textbf{Literature.} The search was limited (a modest number of web searches, Crossref verification of every entry, primary sources read where stated in the bibliography notes). A closer precedent may exist.
\item \textbf{Scope of the application.} Single-threaded tables, names of at most $255$ bytes, one name space per table; no hardware was used.
\end{itemize}

%%%%%%%%%%%%%%%%%%%%%%%%%%%%%%%%%%%%%%%%%%%%%%%%%%%%%%%%%%%%%%%%%%%%%%%%
\section{Conclusion}\label{sec:conclusion}

The last-in-first-out order of scope exit is not just convenient for a scoped symbol table; it is the exact condition under which deletion by clearing is correct, it makes the table's state a function of the live sequence alone, and it turns the analysis of a table with deletions into the classical insertion-only analysis. Using this we gave an exact refusal law for probe-limited tables, tighter probe limits than a Chernoff-type argument, a precise statement of what a key must be and how long it may live, and a ledger showing where a compact layout's memory saving actually comes from. Natural next steps are machine-checked proofs of Theorems~\ref{thm:undo}--\ref{thm:law}, the exact law for other probing rules, a trajectory-level statement about refusals in a run, and measurements of the timing channel on a device.

\section*{Declarations}
\paragraph{Use of AI assistance.} The software, tests, benchmarks and the earlier empirical manuscript that this article extends were produced by an AI coding assistant working in the environment of the first author, under the authors' direction, as recorded in the repository (\code{docs/CONTRIBUTIONS.md}). This article was developed with the same kind of tool: Anthropic's Claude (model Sonnet 5.5, through the Claude Code command-line tool) was used, at one author's direction, to review the earlier work adversarially, to propose and derive the theorems and proofs of Sections~\ref{sec:undo}--\ref{sec:ledger}, to write the verification programs in \code{theory/verify/}, and to draft the text. Every number in the article is produced by those programs, which can be re-run (Appendix~\ref{app:repro}). The tool is not an author; the listed authors are responsible for the content.
\paragraph{Data and code availability.} All sources, programs, traces and raw results are in the public repository~\cite{strata}, directory \code{theory/} for this article.

\bibliographystyle{plainnat}
\bibliography{refs}

\appendix
\section{Proof of Theorem~\ref{thm:law}}\label{app:proof}

Throughout, $k=m-n\ge1$; the case $k=1$ is direct (exactly one cell is empty and it is uniformly located, so the $L$ cells from a uniform start are occupied iff the start lies in the $(m-L)^+$ cells that precede the empty cell by at least $L$). Let $X_j$ be the number of keys with home $j$; $(X_j)$ is multinomial with $\Prob(X=x)=\frac{n!}{m^n}\prod_j\frac1{x_j!}$ on $\{\sum x_j=n\}$.

\paragraph{Step 1: blocks.} Since $n<m$ there are exactly $k$ empty cells and the occupied set depends only on $x$, not on the insertion order~\cite{konheim1966}. Cut the ring at the empty cells: a \emph{block} is an empty cell together with the occupied cells that follow it up to, but excluding, the next empty cell; a block of $g$ cells has one empty cell and $g-1$ occupied ones. No key crosses an empty cell (it would stop there), so the keys homed in a block occupy exactly the cells of that block, $g-1$ of them. Writing $x_1=0,x_2,\dots,x_g$ for the home counts of the cells of the block, the cells $2,\dots,g$ are occupied and the next cell is empty iff
\begin{equation}\label{eq:pattern}
y_u:=x_{u+1}\ \ (1\le u\le g-1)\quad\text{satisfy}\quad S_u:=\sum_{t\le u}y_t\ \ge u\ \ (u<g-1),\qquad S_{g-1}=g-1 .
\end{equation}
Conversely, any ring configuration that is a cyclic concatenation of blocks with patterns~\eqref{eq:pattern} has exactly these empty cells. So home-count vectors with $\sum x_j=n$ correspond bijectively to a placement of $k$ blocks around the ring, each with a pattern~\eqref{eq:pattern}.

\paragraph{Step 2: weight of a block.} Append $y_g:=0$ and put $z_t=y_t-1\ge-1$. Condition~\eqref{eq:pattern} says that the walk $z_1+\dots+z_u$ stays $\ge0$ for $u<g$ and first reaches $-1$ at $u=g$. By the cycle lemma, among the $g$ cyclic rotations of any sequence in $\N^g$ with $\sum y=g-1$ exactly one has this property (no such sequence has a nontrivial period, since the sum $-1$ is not divisible by $g/d$ for a period $d<g$). The weight $\prod_t 1/y_t!$ is rotation invariant, and $\sum_{y\in\N^g,\ \sum y=g-1}\prod_t\frac1{y_t!}=\frac{g^{g-1}}{(g-1)!}$ by the multinomial theorem. Hence the total weight of the patterns of a block of $g$ cells is $\frac1g\frac{g^{g-1}}{(g-1)!}=B(g)$. (With the factors $e^{-\theta}$ of an independent Poisson($\theta$) model this is exactly $\mathrm{Borel}(g;\theta)$, the total progeny of a Poisson($\theta$) branching process~\cite{dwass1969}.)

\paragraph{Step 3: the block containing a cell.} The probability of a home-count vector is proportional to $\prod_j1/x_j!$, so the probability weight of a ring configuration is proportional to the product of the block weights. Count configurations in which the block containing cell $0$ has $g$ cells: cell $0$ occupies one of the $g$ positions of that block, and the other $k-1$ blocks follow in order with $m-g$ cells in total, contributing $\sum_{g_2+\dots+g_k=m-g}\prod B(g_i)=Q_{k-1}(m-g)$. By exchangeability, $\sum_g g\,B(g)Q_{k-1}(m-g)=\frac mk Q_k(m)$ (multiply $Q_k(m)=\sum\prod B(g_i)$ by $m=\sum g_i$ and use symmetry). Therefore
\[
\Prob(\text{block of cell }0\text{ has }g\text{ cells})=\frac{k}{m}\,g\,\pi_g,\qquad \pi_g=\frac{B(g)Q_{k-1}(m-g)}{Q_k(m)},
\]
and, given $g$, the position of cell $0$ in the block is uniform on $\{1,\dots,g\}$.

\paragraph{Step 4: the distance to the next empty cell.} Let $D$ be the number of steps from cell $0$ to the first empty cell at or after it ($D=0$ if cell $0$ is empty). For position $o\in\{2,\dots,g\}$ in a block of $g$ cells, the next empty cell is the first cell of the next block, at distance $g+1-o$, so $D\ge L$ iff $o\le g+1-L$, which holds for $\max(g-L,0)$ positions. Position $1$ has $D=0<L$. Thus
\[
\Prob(D\ge L)=\sum_g\frac kmg\,\pi_g\cdot\frac{(g-L)^+}{g}=\frac km\sum_{g>L}(g-L)\pi_g .
\]
A new key with an independent uniform home sees the first $L$ cells all occupied iff $D\ge L$ from its home; by rotation invariance this is $\Prob(D\ge L)$ at cell $0$. The range $g\le n+1$ holds because a block has at most $n$ occupied cells. $\square$

\section{Reproducing every number}\label{app:repro}\sloppy
Requirements: \code{g++} with C++14, \code{python3} (standard library only), \code{pdflatex} with \code{latexmk} and \code{pgfplots}. From the repository root,
\begin{center}\code{bash theory/verify/run\_all.sh}\end{center}
runs, in order: \code{exact\_law.py} (law against brute-force enumeration in exact rationals), \code{lifo\_check.py} (Theorems~\ref{thm:undo}--\ref{thm:exact} and Proposition~\ref{prop:orphans} on five insertion rules, exhaustive over all home assignments and deletion orders for $m\le8$ as stated in Remark~\ref{rem:general}), \code{design\_table.py} (law against the stored simulation), \code{law\_vs\_table.cpp} with \code{run\_law\_check.py} (law against the real table), \code{design\_rule.py}, \code{proof\_checks.py} (the ledger identities at \LedgerInstants{} instants of the eight traces, the limit and tail-asymptotic corollaries, and the orphan mean: measured \OrphanMeasured{} against \OrphanTheory{} at load $3/4$), \code{space\_model.py} (ledger and the check against the measured minima), \code{adaptive\_key\_demo.cpp}, and \code{make\_figdata.py}, then builds this PDF. All tables, figure data and the numeric macros of the text are written into \code{theory/paper/tables} and \code{theory/paper/data} by these programs. The traces and the measured minima are those of the repository (\code{workloads/traces}, \code{results/raw/20261013-143828-paper}). Approximate run time: eight minutes.

The stack-discipline checks of \code{lifo\_check.py} produced \UndoComparisons{} state comparisons for Theorem~\ref{thm:undo} with \UndoFailures{} failures, the exhaustive check of Theorem~\ref{thm:exact} (\ExactnessCases{} settings: every assignment of homes and every deletion order, linear probing with $m\in\{4,5,6\}$ and all $n\le m$, quadratic probing with $m\in\{4,8\}$ and $n\le4$) \ExactnessFailures{} failures, and the orphan identity of Proposition~\ref{prop:orphans} (all assignments for four small $(m,n)$) \OrphanFailures{} failures.

\end{document}
